%; whizzy paragraph -pdf xpdf -latex ./whizzypdfptex.sh
%; whizzy-paragraph "^\\\\begin{frame}"
% latex beamer presentation.
% platex, latex-beamer でコンパイルすることを想定。 

%     Tokyo Debian Meeting resources
%     Copyright (C) 2009 Junichi Uekawa
%     Copyright (C) 2009 Nobuhiro Iwamatsu

%     This program is free software; you can redistribute it and/or modify
%     it under the terms of the GNU General Public License as published by
%     the Free Software Foundation; either version 2 of the License, or
%     (at your option) any later version.

%     This program is distributed in the hope that it will be useful,
%     but WITHOUT ANY WARRANTY; without even the implied warreanty of
%     MERCHANTABILITY or FITNESS FOR A PARTICULAR PURPOSE.  See the
%     GNU General Public License for more details.

%     You should have received a copy of the GNU General Public License
%     along with this program; if not, write to the Free Software
%     Foundation, Inc., 51 Franklin St, Fifth Floor, Boston, MA  02110-1301 USA

\documentclass[cjk,dvipdfmx,12pt]{beamer}
\usetheme{Tokyo}
\usepackage{monthlypresentation}

%  preview (shell-command (concat "evince " (replace-regexp-in-string
%  "tex$" "pdf"(buffer-file-name)) "&")) 
%  presentation (shell-command (concat "xpdf -fullscreen " (replace-regexp-in-string "tex$" "pdf"(buffer-file-name)) "&"))
%  presentation (shell-command (concat "evince " (replace-regexp-in-string "tex$" "pdf"(buffer-file-name)) "&"))

%http://www.naney.org/diki/dk/hyperref.html
%日本語EUC系環境の時
\AtBeginDvi{\special{pdf:tounicode EUC-UCS2}}
%シフトJIS系環境の時
%\AtBeginDvi{\special{pdf:tounicode 90ms-RKSJ-UCS2}}

\title{第142回東京エリアDebian勉強会 \\　\\Debianでlxcをセットアップしてみよう}
\subtitle{}
\author{Norimitsu Sugimoto (杉本 典充) \\dictoss@live.jp}
\date{2016-08-20}
\logo{\includegraphics[width=8cm]{image200607/openlogo-light.eps}}

\begin{document}

\frame{\titlepage{}}

%\section{}

\begin{frame}{アジェンダ}
  \begin{itemize}
  \item 自己紹介
  \item 仮想化技術について
  \item lxcとは
  \item lxcのインストール
  \item lxcによるコンテナ環境の操作
  \item lxcを実用する
  \item おわりに
  \item 参考資料
  \end{itemize}
\end{frame}


\begin{frame}{自己紹介}
  \begin{itemize}
  \item Norimitsu Sugimoto (杉本 典充)
  \item dictoss@live.jp
  \item Twitter: @dictoss
  \item Debian使って10年以上、FreeBSD使って5年以上
  \item Debian GNU/kFreeBSDが気になっておりウォッチ中
  \item tilegxに少し手を出したがなかなか手ごわい
  \item 仕事はソフトウェア開発者をやってます
  \end{itemize}
\end{frame}


\emtext{仮想化技術について}

\begin{frame}[containsverbatim]{仮想化技術の分類}
  \begin{itemize}
  \item コンテナ型仮想化
  \item 準仮想化型
  \item 完全仮想化型 (エミュレーション型)
  \item 完全仮想化型 (ハイパーバイザ型)
  \end{itemize}
\end{frame}

\begin{frame}[containsverbatim]{仮想化技術のメリット・デメリット}
  \begin{itemize}
  \item 準仮想化型、完全仮想化型
    \begin{itemize}
    \item 物理マシンをエミュレートした仮想マシンとして動作する
    \item 仮想マシン上でもカーネルを動作させる
    \item 物理マシンで動かしていたプログラムはほぼそのまま動く
    \item CPU、メモリ、ディスクを多く消費する
    \end{itemize}
  \item コンテナ型仮想化
    \begin{itemize}
    \item ゲスト環境の動作にカーネルは不要(ホスト環境のカーネル上で動作)
    \item ゲスト環境は、ホスト環境から見るとプロセスとして扱われる
    \item ゲスト環境が利用できるリソースに制約がつく場合がある
    \end{itemize}
  \end{itemize}
\end{frame}

\begin{frame}[containsverbatim]{chroot}
  \begin{itemize}
  \item chrootシステムコールとchrootコマンド
  \item 1982年にビル・ジョイが開発したとされている
  \item \# chroot rootfsdir でコンテナ環境に入れることができる
  \end{itemize}
\end{frame}


\emtext{lxc}

\begin{frame}[containsverbatim]{lxcとは}
  \begin{itemize}
  \item LinuX Containersのことで、省略してlxcと読んでいる
  \item あるディレクトリ配下に実行ファイル、ライブラリ、設定ファイルを適切に配置したrootfsを準備する
  \item rootfsをchroot環境で起動し、仮想マシンのように動かすことができる
  \end{itemize}
\end{frame}


\emtext{lxcのインストール}

\begin{frame}[containsverbatim]{lxcのインストール}
  \begin{itemize}
  \item apt-get install lxc bridge-utils libvirt-bin
  \item lxc-*コマンドが使えるようになります
  \item 今回はlibvirtdでNATネットワークをlxcゲスト環境に割り当てます
  \item cgroupsのmount設定
  \item lxc-checkconfigで動作環境チェック
  \end{itemize}
\end{frame}


\emtext{lxcを使う}

\begin{frame}[containsverbatim]{lxc-create(1)}
  \begin{itemize}
  \item LANG=C SUITE=jessie MIRROR=http://ftp.jp.debian.org/debian lxc-create -n debstudy1 -t debian
  \item 中でdebootstrapを実行してrootfsをダウンロードする
  \item できたlxcのゲスト環境のディレクトリの中
    \begin{itemize}
    \item config
    \item fstab
    \item rootfs
    \end{itemize}
  \end{itemize}
\end{frame}

\begin{frame}[containsverbatim]{lxc-create(2)}
  \begin{itemize}
  \item configを修正して、ネットワークの設定を行う
  \item lxc.network.type = veth
  \item lxc.network.flags = up
  \item lxc.network.link = virbr0
  \item lxc.network.name = eth0
  \item lxc.network.ipv4 = 192.168.122.203/24
  \item lxc.network.ipv4.gateway = 192.168.122.1
  \end{itemize}
\end{frame}

\begin{frame}[containsverbatim]{lxc-ls}
  \begin{itemize}
  \item lxc-lsで/var/lib/lxc配下にいるコンテナ環境の一覧を表示します
  \item --active, --running、--stoppedなどのオプションを指定すると、その状態に絞って一覧を表示できます
  \end{itemize}
\end{frame}

\begin{frame}[containsverbatim]{lxc-destroy}
  \begin{itemize}
  \item lxc-destroy -n debstudy1
  \item コンテナ環境を削除できます
  \item コンテナが実行中の場合は削除できないのがデフォルト
  \item --forceオプションを指定すれば、コンテナが実行中ならstopさせて削除する
  \end{itemize}
\end{frame}

\begin{frame}[containsverbatim]{lxc-start}
  \begin{itemize}
  \item lxc-start -n debstudy1 または lxc-start -n debstudy1 -d
  \item -dまたは--daemonオプションを指定するとバックグラウンドでlxcのゲスト環境が起動する
  \item -dまたは--daemonを指定しない場合はそのターミナル上でlxcのゲスト環境が起動し、コンソールに接続します
  \end{itemize}
\end{frame}

\begin{frame}[containsverbatim]{lxc-stop}
  \begin{itemize}
  \item lxc-stop -n debstudy1
  \item コンテナ環境を終了するよう指示を出します
  \item コンテナ環境の終了とは、initプログラムを終了することをいいます
  \item コンテナ環境でshutdown命令は実行できません
  \end{itemize}
\end{frame}

\begin{frame}[containsverbatim]{lxc-console}
  \begin{itemize}
  \item lxc-console -n debstudy1
  \item lxcのゲスト環境のコンソールに接続します
  \item コンソールを抜ける場合は、「Ctrl+a q」の順に入力してください
  \end{itemize}
\end{frame}


\emtext{lxcを実用する}

\begin{frame}[containsverbatim]{コンテナ環境のセットアップの流れ}
  \begin{itemize}
  \item lxc-createする
  \item lxcのゲスト環境のconfigを書き換えてネットワークを設定
  \item chrootでゲスト環境に入る
    \begin{itemize}
    \item passwdでrootパスワード変更
    \item adduser
    \item apt-get install sudo vim-tiny
    \item visudo
    \end{itemize}
  \item lxc-start
  \item sshログインしてお好みに設定します
  \end{itemize}
\end{frame}

\begin{frame}[containsverbatim]{何にlxcを使うか}
  \begin{itemize}
  \item 一時的な検証で、ホスト環境にいろいろインストールしたくない場合
  \item python2系とpython3系のwsgiアプリを1つのホストで動かしたい場合
  \item ホスト環境はsystemd、ゲスト環境はsysvinitと使い分ける
  \item アプリケーションのクリーンビルドやクリーンインストールをテストする場合
  \item ホスト環境と異なるCPUアーキテクチャのクロス環境がほしい場合
  \end{itemize}
\end{frame}


\begin{frame}[containsverbatim]{おわりに}
  \begin{itemize}
  \item Debian上でlxcを試してみました
  \item 発展系であるLXDやdockerへつなげていきましょう
  \item コンテナは便利ですので試してみてください
  \end{itemize}
\end{frame}


\begin{frame}{参考情報}
  \begin{itemize}
  \item 「LXC」 \url{https://linuxcontainers.org/}
  \item 「LXC - Debian Wiki」 \url{https://wiki.debian.org/LXC}
  \item 「debootstrap を有効活用してみよう」 \url{http://tokyodebian.alioth.debian.org/pdf/debianmeetingresume201304.pdf}
  \end{itemize}
\end{frame}

\end{document}


;;; Local Variables: ***
;;; outline-regexp: "\\([ 	]*\\\\\\(documentstyle\\|documentclass\\|emtext\\|section\\|begin{frame}\\)\\*?[ 	]*[[{]\\|[]+\\)" ***
;;; End: ***
